\documentclass[11pt,a4paper]{article}
\usepackage[top=2cm, bottom=2.3cm, left=1.8cm, right=1.8cm, headsep=0.8cm]{geometry}
\usepackage{amsmath,amssymb,amsfonts}
\usepackage{enumitem}
\usepackage{fancyhdr}
\usepackage[table]{xcolor}
\usepackage{tcolorbox}
\tcbuselibrary{breakable,skins}
\usepackage{tabularx}
\usepackage{multicol}
\usepackage{tikz}
\usetikzlibrary{arrows.meta,patterns,calc}

% High-Contrast Modern Color Palette
\definecolor{primaryDark}{RGB}{12, 32, 84}       % Midnight Navy (Deep)
\definecolor{primaryBlue}{RGB}{18, 72, 160}      % Solid Royal Blue
\definecolor{goldAccent}{RGB}{255, 215, 0}       % Gold highlight
\definecolor{l1Green}{RGB}{14, 120, 62}          % Vivid Forest Green (Level 1)
\definecolor{l2Amber}{RGB}{205, 80, 0}           % Rich Amber / Orange (Level 2)
\definecolor{l3Crimson}{RGB}{180, 20, 20}        % Deep Crimson Red (Level 3)
\definecolor{formulaBg}{RGB}{255, 252, 235}      % Warm Cream-Gold (Formula Callout)
\definecolor{formulaBorder}{RGB}{220, 160, 30}   % Gold Border for Formulas
\definecolor{cardBg}{RGB}{250, 252, 255}         % Clean Crisp Card Background
\definecolor{cardBorder}{RGB}{130, 170, 220}     % Sharp Steel Blue Border
\definecolor{tableStripe}{RGB}{242, 246, 252}    % Alternating Table Row
\definecolor{tableDarkHeader}{RGB}{12, 32, 84}   % Dark Table Header
\definecolor{darkCharcoal}{RGB}{25, 25, 25}      % High-contrast body text

% Page styling (Guaranteed Collision-Free Header)
\pagestyle{fancy}
\fancyhf{}
\setlength{\headheight}{16pt}
\fancyhead[L]{\small\textsf{\textbf{\textcolor{primaryDark}{StudyFAM JEE Revision}} $\vert$ \textbf{\textcolor{primaryBlue}{Day 02: Motion in 2D \& Circular Motion}}}}
\fancyhead[R]{\small\textsf{\textbf{\textcolor{l3Crimson}{Study Target: 8 Hours}}}}
\fancyfoot[C]{\textbf{\textsf{\thepage}}}
\renewcommand{\headrulewidth}{0.8pt}
\renewcommand{\footrulewidth}{0.4pt}

% Custom High-Contrast tcolorboxes
\tcbset{
  dayheader/.style={
    colback=primaryDark,
    coltext=white,
    boxrule=0pt,
    arc=3mm,
    left=14pt, right=14pt, top=10pt, bottom=10pt
  },
  sectionbanner/.style={
    colback=primaryBlue,
    coltext=white,
    boxrule=0pt,
    arc=2mm,
    left=10pt, right=10pt, top=6pt, bottom=6pt
  },
  chapterbanner/.style={
    colback=primaryDark,
    coltext=white,
    boxrule=0pt,
    arc=2mm,
    left=10pt, right=10pt, top=7pt, bottom=7pt
  },
  l1banner/.style={
    colback=l1Green,
    coltext=white,
    boxrule=0pt,
    arc=2mm,
    left=10pt, right=10pt, top=6pt, bottom=6pt
  },
  l2banner/.style={
    colback=l2Amber,
    coltext=white,
    boxrule=0pt,
    arc=2mm,
    left=10pt, right=10pt, top=6pt, bottom=6pt
  },
  l3banner/.style={
    colback=l3Crimson,
    coltext=white,
    boxrule=0pt,
    arc=2mm,
    left=10pt, right=10pt, top=6pt, bottom=6pt
  },
  formulacard/.style={
    enhanced,
    breakable,
    colback=formulaBg,
    colframe=formulaBorder,
    boxrule=1.2pt,
    arc=2mm,
    left=10pt, right=10pt, top=8pt, bottom=8pt
  },
  theorycard/.style={
    enhanced,
    breakable,
    colback=cardBg,
    colframe=cardBorder,
    boxrule=1.2pt,
    arc=2.5mm,
    left=12pt, right=12pt, top=10pt, bottom=10pt
  }
}

\begin{document}

% ==================== DAY HEADER BANNER ====================
\begin{tcolorbox}[dayheader]
\begin{center}
  {\Large \textbf{\textsf{\textcolor{goldAccent}{STUDYFAM JEE INTENSIVE REVISION SERIES}}}}\\[3pt]
  {\LARGE \textbf{\textsf{DAY 02: MOTION IN TWO DIMENSIONS \& CIRCULAR MOTION}}}\\[5pt]
  {\normalsize \textbf{\textsf{Core Topics:}} Projectile Trajectory $\vert$ Incline Projections $\vert$ Circular Kinematics $\vert$ Banking $\vert$ Relative Motion (2D)}\\[3pt]
  {\footnotesize \textbf{Daily Target Time:} 8 Hours \quad $\vert$ \quad \textbf{Curated Questions:} 75 Questions (25 Level-1 + 25 Level-2 + 25 Level-3)}
\end{center}
\end{tcolorbox}

\vspace{6pt}

% ==================== PART 1: TODAY'S PLAN ====================
\begin{tcolorbox}[sectionbanner]
  {\large \textbf{\textsf{PART 1: TODAY'S 8-HOUR ACTION SCHEDULE}}}
\end{tcolorbox}
\vspace{2pt}

\begin{center}
\renewcommand{\arraystretch}{1.25}
\small
\begin{tabularx}{\textwidth}{|c|c|X|c|}
\hline
\rowcolor{tableDarkHeader}
\textcolor{white}{\textbf{Session}} & \textcolor{white}{\textbf{Time}} & \textcolor{white}{\textbf{Focus Area \& High-Yield Strategy}} & \textcolor{white}{\textbf{Duration}} \\
\hline
\rowcolor{white}
\textbf{Block 1} & 09:00 -- 11:00 & \textbf{Theory \& Formula Compendium:} Study the unabridged theory section below. Review 2D independence, trajectory relations, inclined projections, circular acceleration decomposition, and road banking. & 2.0 Hours \\
\hline
\rowcolor{tableStripe}
\textbf{Block 2} & 11:15 -- 12:45 & \textbf{Level-1 Foundation \& Speed Drill (Q1--Q25):} Direct formula application, range symmetry, trajectory equations, uniform circular motion. Target: $\le 1.5$ min/Q. & 1.5 Hours \\
\hline
\rowcolor{white}
-- & 12:45 -- 01:45 & \textit{Lunch Break, Physical Movement \& Mental Recovery} & 1.0 Hour \\
\hline
\rowcolor{tableStripe}
\textbf{Block 3} & 01:45 -- 04:15 & \textbf{Level-2 Core JEE Main Mastery (Q26--Q50):} River-swimmer, rain-man, non-uniform circular motion, tangential/centripetal forces, staircase problems. Target: $\le 2.5$ min/Q. & 2.5 Hours \\
\hline
\rowcolor{white}
\textbf{Block 4} & 04:30 -- 06:15 & \textbf{Level-3 Advanced \& Numerical Challenge (Q51--Q75):} Inclined plane projections, turntable slipping, 2D pursuit, conical pendulum, integer numericals. & 1.75 Hours \\
\hline
\rowcolor{tableStripe}
\textbf{Block 5} & 06:30 -- 07:15 & \textbf{Master Answer Key Audit \& Error Logging:} Check answers, re-derive any missed questions, log tricky traps in your revision notebook. & 0.75 Hour \\
\hline
\rowcolor{tableDarkHeader}
\multicolumn{3}{|r|}{\textcolor{white}{\textbf{TOTAL DEDICATED STUDY TIME}}} & \textcolor{white}{\textbf{8.0 Hours}} \\
\hline
\end{tabularx}
\end{center}

\vspace{5pt}

% ==================== PART 2: OUTCOMES OF THE PLAN ====================
\begin{tcolorbox}[sectionbanner]
  {\large \textbf{\textsf{PART 2: TARGETED LEARNING OUTCOMES}}}
\end{tcolorbox}
\vspace{3pt}

\begin{itemize}[leftmargin=18pt,itemsep=2.5pt,topsep=1pt]
  \item[$\square$] \textbf{\textcolor{primaryDark}{Outcome 1 (Projectile Trajectory \& Geometric Rules):}} Fluently use $y = x\tan\theta\left(1 - \frac{x}{R}\right)$ and $\tan\theta = \frac{4H}{R}$ to solve barrier-clearing and coordinate problems without calculating time.
  \item[$\square$] \textbf{\textcolor{primaryDark}{Outcome 2 (Complementary Angles \& Two-Time Theorems):}} Apply $R(\theta) = R(90^\circ - \theta)$, $R = 4\sqrt{H_1 H_2}$, $t_1 t_2 = \frac{2R}{g}$, and $h = \frac{1}{2}gt_1 t_2$ for crossing identical heights.
  \item[$\square$] \textbf{\textcolor{primaryDark}{Outcome 3 (Projections on Inclined Planes):}} Resolve gravity along and perpendicular to the incline ($g\sin\beta, g\cos\beta$) to find time of flight, range, and perpendicular impact condition $\cot\alpha = 2\tan\beta$.
  \item[$\square$] \textbf{\textcolor{primaryDark}{Outcome 4 (Circular Kinematics \& Vector Nature):}} Distinguish between infinitesimal vector displacement $d\vec{\theta}$ and finite scalar rotation $\theta$, and compute net acceleration $a = \sqrt{a_t^2 + a_c^2}$.
  \item[$\square$] \textbf{\textcolor{primaryDark}{Outcome 5 (Road Banking \& Conical Systems):}} Calculate optimum banking speed $v_0 = \sqrt{Rg\tan\theta}$, friction limits $v_{\max} = \sqrt{Rg\tan(\theta + \phi)}$, and conical pendulum tension $T = \frac{mg}{\cos\theta}$.
  \item[$\square$] \textbf{\textcolor{primaryDark}{Outcome 6 (Relative Motion in 2D):}} Master river-crossing for minimum time ($t_{\min} = d/v_s$) and minimum drift ($x_{\min}$), and determine rain umbrella angles using $\vec{v}_{r/m} = \vec{v}_r - \vec{v}_m$.
  \item[$\square$] \textbf{\textcolor{primaryDark}{Outcome 7 (Relative Angular Velocity):}} Evaluate line-of-sight rotation using $\omega_{A/B} = \frac{(v_{A/B})_\perp}{r_{AB}}$ for intersecting and pursuit paths.
\end{itemize}

\vspace{5pt}

% ==================== PART 3: COMPLETE UNABRIDGED THEORY COMPENDIUM ====================
\begin{tcolorbox}[sectionbanner]
  {\large \textbf{\textsf{PART 3: UNABRIDGED THEORY \& FORMULA REVISION (FROM NOTES)}}}
\end{tcolorbox}
\vspace{4pt}

% -------------------------------------------------------------
% SUBSECTION 1: MOTION IN A PLANE (PROJECTILE MOTION)
% -------------------------------------------------------------
\begin{tcolorbox}[chapterbanner]
  {\Large \textbf{\textsf{1. Motion in a Plane: Projectile Motion (Complete Notes)}}}
\end{tcolorbox}
\vspace{4pt}

\noindent{\large\textbf{\textsf{\textcolor{primaryDark}{1.1 Kinematic Resolution of 2D Motion}}}}\\[3pt]
In two-dimensional motion with constant acceleration $\vec{a} = a_x\hat{i} + a_y\hat{j}$, the horizontal and vertical motions are completely independent:
\begin{itemize}[leftmargin=14pt,itemsep=2pt]
  \item \textbf{Horizontal Component (Uniform Motion):} $u_x = u\cos\theta$, $a_x = 0$, $v_x = u\cos\theta = \text{constant}$, $x = (u\cos\theta)t$.
  \item \textbf{Vertical Component (Uniformly Accelerated Motion):} $u_y = u\sin\theta$, $a_y = -g$, $v_y = u\sin\theta - gt$, $y = (u\sin\theta)t - \frac{1}{2}gt^2$.
  \item \textbf{Instantaneous Velocity Vector:}
  \[
    \vec{v} = v_x\hat{i} + v_y\hat{j} = (u\cos\theta)\hat{i} + (u\sin\theta - gt)\hat{j}, \quad |\vec{v}| = \sqrt{v_x^2 + v_y^2} = \sqrt{u^2 + g^2 t^2 - 2ugt\sin\theta}
  \]
  \item \textbf{Direction of Motion (Angle $\alpha$ with Horizontal):}
  \[
    \tan\alpha = \frac{v_y}{v_x} = \frac{u\sin\theta - gt}{u\cos\theta} = \tan\theta - \frac{gt}{u\cos\theta}
  \]
  \item \textbf{At the Highest Point:} Vertical velocity vanishes ($v_y = 0$), so net speed is minimum but non-zero: $v_{\min} = v_x = u\cos\theta$, directed purely horizontally. Net acceleration is still $g$ downwards ($\vec{a} \perp \vec{v}$).
\end{itemize}

\begin{tcolorbox}[formulacard]
  \textbf{\textcolor{primaryDark}{1.2 Ground-to-Ground Standard Projectile Parameters:}}
  \begin{itemize}[leftmargin=14pt,itemsep=2.5pt]
    \item \textbf{Time of Flight ($T$):} Time required to return to the projection level ($y = 0$):
    \[
      T = \frac{2u_y}{g} = \frac{2u\sin\theta}{g}, \quad t_{\text{ascent}} = t_{\text{descent}} = \frac{u\sin\theta}{g} = \frac{T}{2}
    \]
    \item \textbf{Maximum Height Attained ($H$):} Highest vertical displacement above projection point:
    \[
      H = \frac{u_y^2}{2g} = \frac{u^2\sin^2\theta}{2g} = \frac{1}{8}gT^2 \implies T = \sqrt{\frac{8H}{g}}
    \]
    \item \textbf{Horizontal Range ($R$):} Total horizontal distance traveled during the time of flight:
    \[
      R = u_x T = \frac{2u_x u_y}{g} = \frac{u^2\sin 2\theta}{g}
    \]
    \item \textbf{Fundamental Geometric Relation:}
    \[
      \frac{H}{R} = \frac{u^2\sin^2\theta / 2g}{u^2\sin 2\theta / g} = \frac{\sin^2\theta}{4\sin\theta\cos\theta} = \frac{1}{4}\tan\theta \implies \tan\theta = \frac{4H}{R}
    \]
    \item \textbf{Complementary Angles Property:} For a given initial speed $u$, the range is identical for angles $\theta$ and $(90^\circ - \theta)$:
    \[
      R(\theta) = R(90^\circ - \theta), \quad \frac{H_1}{H_2} = \tan^2\theta, \quad R = 4\sqrt{H_1 H_2}, \quad t_1 t_2 = \frac{2R}{g} \propto R
    \]
    \item \textbf{Maximum Range Condition:} $R$ is maximum at $\theta = 45^\circ$, where $R_{\max} = \dfrac{u^2}{g}$ and $H = \dfrac{R_{\max}}{4}$.
    \item \textbf{Maximum Throw Distance vs. Height:} If an athlete can throw a ball to maximum horizontal distance $X = \frac{u^2}{g}$, the maximum vertical height to which he can throw the same ball is $H_{\max} = \frac{X}{2} = \frac{u^2}{2g}$.
  \end{itemize}
\end{tcolorbox}

\vspace{5pt}
\noindent{\large\textbf{\textsf{\textcolor{primaryDark}{1.3 Equation of Trajectory (Path of Projectile)}}}}\\[3pt]
Eliminating time $t = \frac{x}{u\cos\theta}$ gives the parabolic trajectory:
\[
  y = x\tan\theta - \frac{gx^2}{2u^2\cos^2\theta} = x\tan\theta\left(1 - \frac{x}{R}\right)
\]
\textit{Crucial Exam Note:} Any equation of the form $y = Ax - Bx^2$ represents a projectile trajectory with $\tan\theta = A$ and range $R = \frac{A}{B}$.

\vspace{5pt}
\noindent{\large\textbf{\textsf{\textcolor{primaryDark}{1.4 Two-Point Crossing Theorems (From Notes)}}}}\\[3pt]
If a projectile crosses the same vertical height $h$ at two instants $t_1$ and $t_2$, corresponding to horizontal distances $x$ and $y$ from the origin:
\begin{enumerate}[label=(\alph*),leftmargin=14pt,itemsep=2pt]
  \item $x + y = R$ (the sum of the horizontal coordinates equals the total range).
  \item $t_1 + t_2 = T$ (sum of crossing times equals total time of flight).
  \item $h = \frac{1}{2}gt_1 t_2$ (height determined entirely by crossing times).
  \item \textbf{Average Velocity:} The average velocity of the projectile between the two points at equal height is purely horizontal and equal to $u\cos\theta$.
\end{enumerate}

\vspace{5pt}
\noindent{\large\textbf{\textsf{\textcolor{primaryDark}{1.5 Horizontal Projection from a Height $h$}}}}\\[3pt]
When a body is projected horizontally with velocity $u$ from the top of a tower of height $h$:
\begin{itemize}[leftmargin=14pt,itemsep=1.5pt]
  \item $u_x = u, a_x = 0 \implies x = ut$; \quad $u_y = 0, a_y = g \implies y = \frac{1}{2}gt^2$.
  \item \textbf{Time of Flight:} $T = \sqrt{\dfrac{2h}{g}}$ (independent of projection speed $u$!).
  \item \textbf{Horizontal Range:} $R = uT = u\sqrt{\dfrac{2h}{g}}$.
  \item \textbf{Impact Velocity:} $v_x = u$, $v_y = gt = \sqrt{2gh} \implies v = \sqrt{u^2 + 2gh}$, inclined at $\tan\theta = \dfrac{v_y}{v_x} = \dfrac{\sqrt{2gh}}{u}$.
\end{itemize}

\vspace{5pt}
\noindent{\large\textbf{\textsf{\textcolor{primaryDark}{1.6 Angled Projection from a Height $h$ (From Notes)}}}}\\[3pt]
\begin{itemize}[leftmargin=14pt,itemsep=2pt]
  \item \textbf{Case I: Projected Upwards at Angle $\theta$ Above Horizontal:}
  Time taken to strike the ground $t_1$:
  \[
    -h = (u\sin\theta)t_1 - \frac{1}{2}gt_1^2 \implies t_1 = \frac{T + \sqrt{T^2 + 8h/g}}{2}, \quad \text{where } T = \frac{2u\sin\theta}{g}
  \]
  Horizontal range: $R = (u\cos\theta)t_1$.
  \item \textbf{Case II: Projected Downwards at Angle $\theta$ Below Horizontal:}
  Time taken to strike the ground $t_2$:
  \[
    h = (u\sin\theta)t_2 + \frac{1}{2}gt_2^2 \implies t_2 = \frac{-T + \sqrt{T^2 + 8h/g}}{2}
  \]
  Horizontal range: $R = (u\cos\theta)t_2$.
\end{itemize}

\vspace{5pt}
\noindent{\large\textbf{\textsf{\textcolor{primaryDark}{1.7 Projectile Motion on an Inclined Plane (From Notes)}}}}\\[3pt]
Let the inclined plane make an angle $\beta$ with the horizontal, and the projectile be launched at an angle $\alpha$ relative to the inclined surface:
\begin{itemize}[leftmargin=14pt,itemsep=2.5pt]
  \item \textbf{Up the Incline:} Acceleration components: $a_\parallel = -g\sin\beta$, $a_\perp = -g\cos\beta$.
  \[
    T = \frac{2u_\perp}{g_\perp} = \frac{2u\sin\alpha}{g\cos\beta}, \quad H = \frac{u_\perp^2}{2g_\perp} = \frac{u^2\sin^2\alpha}{2g\cos\beta}
  \]
  \[
    R_{\text{up}} = \frac{2u^2\sin\alpha\cos(\alpha + \beta)}{g\cos^2\beta}, \quad R_{\max} = \frac{u^2}{g(1 + \sin\beta)} \quad \left(\text{at } \alpha = \frac{\pi}{4} - \frac{\beta}{2}\right)
  \]
  \item \textbf{Down the Incline:} Acceleration components: $a_\parallel = +g\sin\beta$, $a_\perp = -g\cos\beta$.
  \[
    T = \frac{2u\sin\alpha}{g\cos\beta}, \quad R_{\text{down}} = \frac{2u^2\sin\alpha\cos(\alpha - \beta)}{g\cos^2\beta}, \quad R_{\max} = \frac{u^2}{g(1 - \sin\beta)} \quad \left(\text{at } \alpha = \frac{\pi}{4} + \frac{\beta}{2}\right)
  \]
  \item \textbf{Perpendicular Strike Condition:} If the projectile strikes the inclined plane normally ($v_\parallel = 0$ at landing):
  \[
    \tan\alpha = \frac{1}{2\tan\beta} \iff \cot\alpha = 2\tan\beta
  \]
\end{itemize}

\vspace{10pt}

% -------------------------------------------------------------
% SUBSECTION 2: CIRCULAR MOTION KINEMATICS & DYNAMICS
% -------------------------------------------------------------
\begin{tcolorbox}[chapterbanner]
  {\Large \textbf{\textsf{2. Circular Motion Kinematics \& Dynamics (Complete Notes)}}}
\end{tcolorbox}
\vspace{4pt}

\noindent{\large\textbf{\textsf{\textcolor{primaryDark}{2.1 Angular Variables \& Vector Character (Key Tips From Notes)}}}}\\[3pt]
\begin{itemize}[leftmargin=14pt,itemsep=2pt]
  \item \textbf{Radius Vector ($\vec{r}$):} Vector directed from the center of curvature to the revolving particle. It has constant magnitude $r$, but continuous varying direction ($d\vec{r}/dt \ne 0$).
  \item \textbf{Angular Displacement ($\theta$):} $\theta = \frac{\text{arc length } s}{\text{radius } r}$ (measured in radians).
  \begin{tcolorbox}[formulacard]
    \textbf{\textcolor{l3Crimson}{Fundamental Rule of Angular Displacement:}}
    \begin{itemize}[leftmargin=12pt,itemsep=1.5pt]
      \item \textbf{Infinitesimal angular displacement $d\vec{\theta}$ is a TRUE VECTOR:} It obeys the commutative law of vector addition ($d\vec{\theta}_1 + d\vec{\theta}_2 = d\vec{\theta}_2 + d\vec{\theta}_1$).
      \item \textbf{Finite angular displacement $\theta$ is a SCALAR:} Finite rotations about different axes do not commute ($\vec{\theta}_1 + \vec{\theta}_2 \ne \vec{\theta}_2 + \vec{\theta}_1$).
    \end{itemize}
  \end{tcolorbox}
  \item \textbf{Angular Velocity ($\omega$):}
    \begin{itemize}[leftmargin=12pt,itemsep=1pt]
      \item \textit{Average:} $\omega_{\text{avg}} = \dfrac{\Delta\theta}{\Delta t}$ (scalar quantity).
      \item \textit{Instantaneous:} $\vec{\omega} = \dfrac{d\vec{\theta}}{dt}$ (axial vector along the rotation axis given by right-hand thumb rule). Unit: $\text{rad/s}$.
      \item \textit{Uniform Rotation:} $\omega = \dfrac{2\pi}{T} = 2\pi f = 2\pi n$, where $T = \frac{1}{n}$ is the time period and $n$ is frequency in $\text{rev/s}$ (or $\text{rpm}/60$).
    \end{itemize}
  \item \textbf{Linear vs. Angular Velocity Relation:}
  \[
    v = \omega r \quad \iff \quad \vec{v} = \vec{\omega} \times \vec{r}
  \]
\end{itemize}

\vspace{5pt}
\noindent{\large\textbf{\textsf{\textcolor{primaryDark}{2.2 Acceleration Decomposition in Circular Motion}}}}\\[3pt]
Differentiating velocity vector $\vec{v} = \vec{\omega}\times\vec{r}$ with respect to time:
\[
  \vec{a} = \frac{d\vec{v}}{dt} = \frac{d\vec{\omega}}{dt}\times\vec{r} + \vec{\omega}\times\frac{d\vec{r}}{dt} = \vec{\alpha}\times\vec{r} + \vec{\omega}\times\vec{v} = \vec{a}_t + \vec{a}_c
\]
\begin{center}
\renewcommand{\arraystretch}{1.25}
\small
\begin{tabularx}{\linewidth}{|>{\raggedright\arraybackslash}m{3.2cm}|>{\centering\arraybackslash}m{3.0cm}|>{\centering\arraybackslash}m{3.2cm}|>{\raggedright\arraybackslash}X|}
\hline
\rowcolor{tableDarkHeader}
\textcolor{white}{\textbf{Component}} & \textcolor{white}{\textbf{Formula}} & \textcolor{white}{\textbf{Direction}} & \textcolor{white}{\textbf{Physical Role}} \\
\hline
\rowcolor{white}
\textbf{Tangential Accel.} ($\vec{a}_t$) & $a_t = \dfrac{dv}{dt} = \alpha r$ & Tangent to path & Changes the \textbf{magnitude} of velocity (speed). Zero in uniform circular motion. \\
\hline
\rowcolor{tableStripe}
\textbf{Centripetal Accel.} ($\vec{a}_c$) & $a_c = \dfrac{v^2}{r} = \omega^2 r = \omega v$ & Towards center (radial) & Changes the \textbf{direction} of velocity. Always present in any curved motion! \\
\hline
\rowcolor{white}
\textbf{Net Acceleration} ($\vec{a}$) & $a = \sqrt{a_t^2 + a_c^2}$ & $\tan\phi = \dfrac{a_c}{a_t}$ & $a = \sqrt{\left(\dfrac{dv}{dt}\right)^2 + \left(\dfrac{v^2}{r}\right)^2} = \sqrt{(\alpha r)^2 + (\omega^2 r)^2}$. \\
\hline
\end{tabularx}
\end{center}

\vspace{4pt}
\begin{tcolorbox}[formulacard]
  \textbf{\textcolor{primaryDark}{2.3 Road Banking, Friction Limits \& Bending (From Notes):}}
  \begin{itemize}[leftmargin=14pt,itemsep=2.5pt]
    \item \textbf{Level Unbanked Circular Road:} Necessary centripetal force is provided solely by static friction:
    \[
      f_s \le \mu_s mg \implies \frac{mv^2}{R} \le \mu_s mg \implies v_{\max} = \sqrt{\mu_s R g}
    \]
    \item \textbf{Optimum Banking Speed (Friction-Free Speed):} Road banked at angle $\theta$ with horizontal:
    \[
      N\sin\theta = \frac{mv_0^2}{R}, \quad N\cos\theta = mg \implies \tan\theta = \frac{v_0^2}{Rg} \implies v_0 = \sqrt{Rg\tan\theta}
    \]
    \item \textbf{Banked Road with Friction ($\mu_s$):}
    \[
      v_{\max} = \sqrt{Rg\left(\frac{\mu_s + \tan\theta}{1 - \mu_s\tan\theta}\right)} = \sqrt{Rg\tan(\theta + \phi)}, \quad \text{where } \tan\phi = \mu_s
    \]
    \[
      v_{\min} = \sqrt{Rg\left(\frac{\tan\theta - \mu_s}{1 + \mu_s\tan\theta}\right)} = \sqrt{Rg\tan(\theta - \phi)}
    \]
    \item \textbf{Bending of a Cyclist:} A cyclist negotiating a turn of radius $r$ at speed $v$ bends inward at angle $\theta$ with the vertical:
    \[
      \tan\theta = \frac{v^2}{rg}
    \]
  \end{itemize}
\end{tcolorbox}

\vspace{5pt}
\noindent{\large\textbf{\textsf{\textcolor{primaryDark}{2.4 Relative Angular Velocity (From Notes)}}}}\\[3pt]
The relative angular velocity of particle $A$ with respect to particle $B$ is the rate at which the line of sight $AB$ rotates:
\[
  \omega_{A/B} = \frac{(v_{A/B})_\perp}{r_{AB}} = \frac{\text{Component of relative velocity } \perp \text{ to line } AB}{\text{Separation distance } r_{AB}}
\]
If $\vec{v}_A$ makes angle $\theta_1$ and $\vec{v}_B$ makes angle $\theta_2$ with the line connecting them:
\[
  \omega_{A/B} = \frac{v_A\sin\theta_1 + v_B\sin\theta_2}{r_{AB}}
\]

\vspace{5pt}
\noindent{\large\textbf{\textsf{\textcolor{primaryDark}{2.5 Relative Motion in 2D: River-Swimmer \& Rain-Man}}}}\\[3pt]
\begin{itemize}[leftmargin=14pt,itemsep=2pt]
  \item \textbf{River-Swimmer Problems:} River speed $v_r$, swimmer speed in still water $v_s$, width $d$.
  \begin{itemize}[leftmargin=12pt,itemsep=1pt]
    \item \textit{Condition for Shortest Time:} Swim directly perpendicular to river current ($\theta = 90^\circ$).
    \[
      t_{\min} = \frac{d}{v_s}, \quad \text{Drift along bank: } x = v_r t_{\min} = \frac{v_r d}{v_s}
    \]
    \item \textit{Condition for Shortest Path (Zero Drift):} Possible only if $v_s > v_r$. Swimmer must head upstream at angle $\theta = \sin^{-1}\left(\frac{v_r}{v_s}\right)$ with the normal:
    \[
      t = \frac{d}{v_{\text{eff}}} = \frac{d}{\sqrt{v_s^2 - v_r^2}}
    \]
    \item \textit{Minimum Drift when $v_s < v_r$:} Swimmer heads upstream at angle $\theta = \sin^{-1}\left(\frac{v_s}{v_r}\right)$, giving minimum drift:
    \[
      x_{\min} = d\sqrt{\left(\frac{v_r}{v_s}\right)^2 - 1}
    \]
  \end{itemize}
  \item \textbf{Rain-Man Problems:} Relative velocity of rain with respect to man: $\vec{v}_{r/m} = \vec{v}_r - \vec{v}_m$.
  To protect from rain, hold umbrella along the direction of $\vec{v}_{r/m}$:
  \[
    \tan\phi = \frac{|\vec{v}_m|}{|\vec{v}_r|} \quad \text{(measured from the vertical)}
  \]
\end{itemize}

\newpage

\begin{tcolorbox}[l1banner]
  {\Large \textbf{\textsf{LEVEL 1: FOUNDATION \& SPEED DRILLS (25 Questions)}}}\\[2pt]
  {\footnotesize \textbf{Focus:} Direct Formula Application $\vert$ Range Symmetry $\vert$ Trajectory Basics $\vert$ Target Time: $\le 1.5$ min/Q}
\end{tcolorbox}
\vspace{6pt}

\begin{enumerate}[label=\textbf{\textcolor{l1Green}{Q\arabic*.}},leftmargin=24pt,itemsep=10pt]

  % Q1
  \item Two particles start simultaneously from the same point and move along two straight lines, one with uniform velocity $v$ and the other with uniform acceleration $a$. If $\alpha$ is the angle between their lines of motion, the least value of relative velocity will occur at time:
  \begin{multicols}{2}
  \begin{enumerate}[label=\textbf{\textcolor{primaryDark}{(\Alph*)}},itemsep=2pt]
    \item $\dfrac{v\sin\alpha}{a}$
    \item $\dfrac{v\cos\alpha}{a}$
    \item $\dfrac{v\tan\alpha}{a}$
    \item $\dfrac{v\cot\alpha}{a}$
  \end{enumerate}
  \end{multicols}

  % Q2
  \item In a car race, car $A$ takes time $t_0$ less to finish than car $B$ and passes the finishing point with a velocity $v_0$ more than car $B$. Both cars start from rest and travel with constant accelerations $a_1$ and $a_2$ respectively. The ratio $\dfrac{v_0}{t_0}$ is equal to:
  \begin{multicols}{2}
  \begin{enumerate}[label=\textbf{\textcolor{primaryDark}{(\Alph*)}},itemsep=2pt]
    \item $\dfrac{a_1 + a_2}{2}$
    \item $\dfrac{a_1 a_2}{a_1 + a_2}$
    \item $\sqrt{a_1 a_2}$
    \item $\dfrac{a_2^2}{a_1}$
  \end{enumerate}
  \end{multicols}

  % Q3
  \item A man can swim at a speed of $5\text{ km/h}$ in still water. A $1\text{ km}$ wide river flows at the rate of $3\text{ km/h}$. If the man wishes to swim across the river directly opposite to the starting point, the time he will take to cross the river is:
  \begin{multicols}{2}
  \begin{enumerate}[label=\textbf{\textcolor{primaryDark}{(\Alph*)}},itemsep=2pt]
    \item $10\text{ min}$
    \item $13\text{ min}$
    \item $15\text{ min}$
    \item $18\text{ min}$
  \end{enumerate}
  \end{multicols}

  % Q4
  \item A projectile has the same horizontal range $R$ for two different angles of projection. If $t_1$ and $t_2$ be the times of flight in the two cases, then the product of the two times of flight satisfies:
  \begin{multicols}{2}
  \begin{enumerate}[label=\textbf{\textcolor{primaryDark}{(\Alph*)}},itemsep=2pt]
    \item $t_1 t_2 \propto R$
    \item $t_1 t_2 \propto R^2$
    \item $t_1 t_2 \propto \dfrac{1}{R}$
    \item $t_1 t_2 \propto \dfrac{1}{R^2}$
  \end{enumerate}
  \end{multicols}

  % Q5
  \item If the angles of projection of a projectile with the same initial velocity exceed or fall short of $45^\circ$ by equal amounts $\alpha$, then the ratio of their horizontal ranges is:
  \begin{multicols}{2}
  \begin{enumerate}[label=\textbf{\textcolor{primaryDark}{(\Alph*)}},itemsep=2pt]
    \item $1 : 2$
    \item $1 : 3$
    \item $1 : 4$
    \item $1 : 1$
  \end{enumerate}
  \end{multicols}

  % Q6
  \item The equation of trajectory of a projectile is given by $y = \sqrt{3}x - \dfrac{gx^2}{2}$. The angle of projection $\theta$ with the horizontal is:
  \begin{multicols}{2}
  \begin{enumerate}[label=\textbf{\textcolor{primaryDark}{(\Alph*)}},itemsep=2pt]
    \item $\tan\theta = \dfrac{1}{\sqrt{3}}$
    \item $\tan\theta = \sqrt{3}$
    \item $\theta = \dfrac{\pi}{2}$
    \item $\theta = 0^\circ$
  \end{enumerate}
  \end{multicols}

  % Q7
  \item An airplane flying horizontally at a height of $1500\text{ m}$ with a velocity of $200\text{ m/s}$ passes directly overhead an anti-aircraft gun. The angle with the horizontal at which the gun should be fired with a muzzle velocity of $400\text{ m/s}$ to hit the plane is:
  \begin{multicols}{2}
  \begin{enumerate}[label=\textbf{\textcolor{primaryDark}{(\Alph*)}},itemsep=2pt]
    \item $90^\circ$
    \item $60^\circ$
    \item $30^\circ$
    \item $45^\circ$
  \end{enumerate}
  \end{multicols}

  % Q8
  \item A body projected at an angle with the horizontal has a range of $300\text{ m}$. If the time of flight is $6\text{ s}$, then the horizontal component of velocity is:
  \begin{multicols}{2}
  \begin{enumerate}[label=\textbf{\textcolor{primaryDark}{(\Alph*)}},itemsep=2pt]
    \item $30\text{ m/s}$
    \item $50\text{ m/s}$
    \item $40\text{ m/s}$
    \item $45\text{ m/s}$
  \end{enumerate}
  \end{multicols}

  % Q9
  \item A boy can throw a stone up to a maximum vertical height of $10\text{ m}$. The maximum horizontal distance that the boy can throw the same stone is:
  \begin{multicols}{2}
  \begin{enumerate}[label=\textbf{\textcolor{primaryDark}{(\Alph*)}},itemsep=2pt]
    \item $10\text{ m}$
    \item $10\sqrt{2}\text{ m}$
    \item $15\text{ m}$
    \item $20\text{ m}$
  \end{enumerate}
  \end{multicols}

  % Q10
  \item A projectile is given an initial velocity of $\vec{v} = (\hat{i} + 2\hat{j})\text{ m/s}$, where $\hat{i}$ is along the horizontal ground and $\hat{j}$ is vertically upwards. If $g = 10\text{ m/s}^2$, the equation of its trajectory is:
  \begin{multicols}{2}
  \begin{enumerate}[label=\textbf{\textcolor{primaryDark}{(\Alph*)}},itemsep=2pt]
    \item $y = x - 5x^2$
    \item $y = 2x - 5x^2$
    \item $4y = 2x - 5x^2$
    \item $4y = 2x - 25x^2$
  \end{enumerate}
  \end{multicols}

  % Q11
  \item The equation of trajectory of a projectile is given by $y = 4x - \dfrac{5}{39}x^2$. The horizontal range of the projectile is:
  \begin{multicols}{2}
  \begin{enumerate}[label=\textbf{\textcolor{primaryDark}{(\Alph*)}},itemsep=2pt]
    \item $31.2\text{ m}$
    \item $15.6\text{ m}$
    \item $7.8\text{ m}$
    \item $23.4\text{ m}$
  \end{enumerate}
  \end{multicols}

  % Q12
  \item A projectile is fired from the surface of the earth with a velocity of $5\text{ m/s}$ at angle $\theta$ with the horizontal. Another projectile fired from another planet with a velocity of $3\text{ m/s}$ at the same angle follows an identical trajectory. The acceleration due to gravity on the planet is: (Take $g = 9.8\text{ m/s}^2$)
  \begin{multicols}{2}
  \begin{enumerate}[label=\textbf{\textcolor{primaryDark}{(\Alph*)}},itemsep=2pt]
    \item $3.53\text{ m/s}^2$
    \item $5.88\text{ m/s}^2$
    \item $16.3\text{ m/s}^2$
    \item $2.45\text{ m/s}^2$
  \end{enumerate}
  \end{multicols}

  % Q13
  \item A projectile is thrown at an angle of $40^\circ$ with the horizontal and its range is $R_1$. Another projectile is thrown at an angle of $40^\circ$ with the vertical and its range is $R_2$. The relation between $R_1$ and $R_2$ is:
  \begin{multicols}{2}
  \begin{enumerate}[label=\textbf{\textcolor{primaryDark}{(\Alph*)}},itemsep=2pt]
    \item $R_1 = R_2$
    \item $R_1 = 2R_2$
    \item $2R_1 = R_2$
    \item $R_1 = \dfrac{4}{5}R_2$
  \end{enumerate}
  \end{multicols}

  % Q14
  \item A particle is projected at an angle of elevation $\alpha$. After $t$ seconds, its position vector appears to have an angle of elevation $\beta$ as seen from the point of projection. The initial projection speed $u$ is:
  \begin{multicols}{2}
  \begin{enumerate}[label=\textbf{\textcolor{primaryDark}{(\Alph*)}},itemsep=2pt]
    \item $\dfrac{gt\cos\beta}{2\sin(\alpha - \beta)}$
    \item $\dfrac{gt\sin\beta}{2\sin(\alpha - \beta)}$
    \item $\dfrac{gt}{2\sin(\alpha - \beta)}$
    \item $\dfrac{gt\cos\beta}{\sin(\alpha - \beta)}$
  \end{enumerate}
  \end{multicols}

  % Q15
  \item A particle is projected with speed $u$ from the ground at an angle $\theta$ with the horizontal. The radius of curvature of the trajectory at the point where its velocity vector is perpendicular to the initial velocity vector is:
  \begin{multicols}{2}
  \begin{enumerate}[label=\textbf{\textcolor{primaryDark}{(\Alph*)}},itemsep=2pt]
    \item $\dfrac{u^2\cos^2\theta}{g}$
    \item $\dfrac{u^2\cot^2\theta}{g\sin\theta}$
    \item $\dfrac{u^2}{g}$
    \item $\dfrac{u^2\tan^2\theta}{g\cos\theta}$
  \end{enumerate}
  \end{multicols}

  % Q16
  \item A body of mass $m$ is projected horizontally with velocity $v$ from the top of a tower of height $h$ and reaches the ground at distance $x$ from the foot. If a second body of mass $2m$ is projected horizontally from height $2h$ and reaches distance $2x$, the horizontal velocity of the second body is:
  \begin{multicols}{2}
  \begin{enumerate}[label=\textbf{\textcolor{primaryDark}{(\Alph*)}},itemsep=2pt]
    \item $v$
    \item $\sqrt{2}v$
    \item $2v$
    \item $\dfrac{v}{2}$
  \end{enumerate}
  \end{multicols}

  % Q17
  \item The velocity of projection of an oblique projectile is $\vec{u} = (6\hat{i} + 8\hat{j})\text{ m/s}$. Taking $g = 10\text{ m/s}^2$, the horizontal range of the projectile is:
  \begin{multicols}{2}
  \begin{enumerate}[label=\textbf{\textcolor{primaryDark}{(\Alph*)}},itemsep=2pt]
    \item $4.9\text{ m}$
    \item $9.6\text{ m}$
    \item $19.6\text{ m}$
    \item $14.4\text{ m}$
  \end{enumerate}
  \end{multicols}

  % Q18
  \item A liquid of mass $M$ fills a narrow tube $AB$ of length $L$. The tube is rotated in a horizontal plane about end $A$ with constant angular velocity $\omega$. The total force exerted by the rotating liquid on the closed end $B$ is:
  \begin{multicols}{2}
  \begin{enumerate}[label=\textbf{\textcolor{primaryDark}{(\Alph*)}},itemsep=2pt]
    \item $\dfrac{1}{2}ML\omega^2$
    \item $ML\omega^2$
    \item $\dfrac{1}{3}ML\omega^2$
    \item $\dfrac{1}{4}ML\omega^2$
  \end{enumerate}
  \end{multicols}

  % Q19
  \item A stone tied to the end of a string of length $1\text{ m}$ is whirled in a horizontal circle with a constant speed. If the stone makes $22$ revolutions in $44\text{ seconds}$, the magnitude and direction of the acceleration of the stone are:
  \begin{multicols}{2}
  \begin{enumerate}[label=\textbf{\textcolor{primaryDark}{(\Alph*)}},itemsep=2pt]
    \item $\pi^2\text{ m/s}^2$ towards the center
    \item $\pi^2\text{ m/s}^2$ away from center
    \item $\pi^2\text{ m/s}^2$ along the tangent
    \item $\dfrac{\pi^2}{4}\text{ m/s}^2$ towards center
  \end{enumerate}
  \end{multicols}

  % Q20
  \item A particle moves along a circular path of radius $R = 100\text{ m}$. Its speed varies with time as $v(t) = 2t^2 + t$ (in $\text{m/s}$). The magnitude of the net acceleration of the particle at $t = 2\text{ s}$ is approximately:
  \begin{multicols}{2}
  \begin{enumerate}[label=\textbf{\textcolor{primaryDark}{(\Alph*)}},itemsep=2pt]
    \item $21\text{ m/s}^2$
    \item $9.06\text{ m/s}^2$
    \item $13.5\text{ m/s}^2$
    \item $10.0\text{ m/s}^2$
  \end{enumerate}
  \end{multicols}

  % Q21
  \item A particle is projected from the origin. The $x$ and $y$ components of its displacement vary with time as $x(t) = 6t$ and $y(t) = 8t - 5t^2$ (in meters). The initial projection speed of the particle is:
  \begin{multicols}{2}
  \begin{enumerate}[label=\textbf{\textcolor{primaryDark}{(\Alph*)}},itemsep=2pt]
    \item $6\text{ m/s}$
    \item $8\text{ m/s}$
    \item $10\text{ m/s}$
    \item $14\text{ m/s}$
  \end{enumerate}
  \end{multicols}

  % Q22
  \item A particle is projected with a velocity $v$ such that its horizontal range is twice the greatest height attained ($R = 2H$). If $g$ is acceleration due to gravity, the range $R$ is:
  \begin{multicols}{2}
  \begin{enumerate}[label=\textbf{\textcolor{primaryDark}{(\Alph*)}},itemsep=2pt]
    \item $\dfrac{4v^2}{5g}$
    \item $\dfrac{5v^2}{4g}$
    \item $\dfrac{4v^2}{3g}$
    \item $\dfrac{2v^2}{5g}$
  \end{enumerate}
  \end{multicols}

  % Q23
  \item A body moves in a circular path of radius $r$ with a constant speed $v$. Which of the following statements is TRUE?
  \begin{enumerate}[label=\textbf{\textcolor{primaryDark}{(\Alph*)}},itemsep=2pt]
    \item Its velocity is constant.
    \item Its acceleration vector is constant.
    \item It possesses an acceleration of constant magnitude directed radially inward.
    \item Its acceleration varies in magnitude with time.
  \end{enumerate}

  % Q24
  \item A particle is moving along a circle of radius $r$ centered at $O$ with a constant speed $v$. What is the magnitude of the change in velocity when it turns through an angle $\angle AOB = 40^\circ$?
  \begin{multicols}{2}
  \begin{enumerate}[label=\textbf{\textcolor{primaryDark}{(\Alph*)}},itemsep=2pt]
    \item $2v\sin 20^\circ$
    \item $4v\sin 40^\circ$
    \item $2v\sin 40^\circ$
    \item $v\sin 20^\circ$
  \end{enumerate}
  \end{multicols}

  % Q25
  \item An aircraft executes a horizontal loop of radius $R = 1\text{ km}$ with a steady speed of $900\text{ km/h}$. The centripetal acceleration of the aircraft is:
  \begin{multicols}{2}
  \begin{enumerate}[label=\textbf{\textcolor{primaryDark}{(\Alph*)}},itemsep=2pt]
    \item $62.5\text{ m/s}^2$
    \item $30.0\text{ m/s}^2$
    \item $32.5\text{ m/s}^2$
    \item $72.5\text{ m/s}^2$
  \end{enumerate}
  \end{multicols}

\end{enumerate}

\newpage

% ==================== LEVEL 2 ====================
\begin{tcolorbox}[l2banner]
  {\Large \textbf{\textsf{LEVEL 2: JEE MAIN CORE MASTERY (25 Questions)}}}\\[2pt]
  {\footnotesize \textbf{Focus:} Two-Step Reasoning $\vert$ River-Swimmer $\vert$ Rain-Man $\vert$ Road Banking $\vert$ Target Time: $\le 2.5$ min/Q}
\end{tcolorbox}
\vspace{6pt}

\begin{enumerate}[label=\textbf{\textcolor{l2Amber}{Q\arabic*.}},leftmargin=24pt,itemsep=10pt,start=26]

  % Q26
  \item A man holds an umbrella at $30^\circ$ with the vertical to keep himself dry. When he runs at a speed of $10\text{ m/s}$, he finds raindrops hitting him vertically. The speeds of raindrops with respect to the running man and with respect to the earth are respectively:
  \begin{multicols}{2}
  \begin{enumerate}[label=\textbf{\textcolor{primaryDark}{(\Alph*)}},itemsep=2pt]
    \item $20\text{ m/s}, 10\text{ m/s}$
    \item $10\text{ m/s}, 20\sqrt{3}\text{ m/s}$
    \item $10\sqrt{3}\text{ m/s}, 20\text{ m/s}$
    \item $20\text{ m/s}, 10\sqrt{3}\text{ m/s}$
  \end{enumerate}
  \end{multicols}

  % Q27
  \item A bird is flying due east with a velocity of $4\text{ m/s}$. The wind starts blowing with a velocity of $3\text{ m/s}$ due north. The relative velocity of the bird with respect to the wind is directed:
  \begin{multicols}{2}
  \begin{enumerate}[label=\textbf{\textcolor{primaryDark}{(\Alph*)}},itemsep=2pt]
    \item $\tan^{-1}(3/4)$ South of East
    \item $\tan^{-1}(4/3)$ South of East
    \item $\tan^{-1}(3/4)$ North of East
    \item $\tan^{-1}(4/3)$ East of South
  \end{enumerate}
  \end{multicols}

  % Q28
  \item A man can row a boat at $4\text{ km/h}$ in still water. He crosses a river of width $4\text{ km}$ where the current flows at $2\text{ km/h}$. In what direction should he head the boat to cross the river in the shortest time, and what is this minimum time?
  \begin{multicols}{2}
  \begin{enumerate}[label=\textbf{\textcolor{primaryDark}{(\Alph*)}},itemsep=2pt]
    \item Parallel to current, $2\text{ hours}$
    \item Perpendicular to current, $2\text{ hours}$
    \item At $60^\circ$ upstream, $1\text{ hour}$
    \item Perpendicular to current, $1\text{ hour}$
  \end{enumerate}
  \end{multicols}

  % Q29
  \item A stationary observer sees rain falling with a velocity of $20\text{ m/s}$ at an angle of $30^\circ$ with the vertical. The velocity with which the observer should move so that rain appears to fall vertically is:
  \begin{multicols}{2}
  \begin{enumerate}[label=\textbf{\textcolor{primaryDark}{(\Alph*)}},itemsep=2pt]
    \item $15\text{ m/s}$
    \item $30\text{ m/s}$
    \item $10\text{ m/s}$
    \item $20\text{ m/s}$
  \end{enumerate}
  \end{multicols}

  % Q30
  \item A particle has an initial velocity $\vec{u} = (3\hat{i} + 4\hat{j})\text{ m/s}$ and a constant acceleration $\vec{a} = (0.4\hat{i} + 0.3\hat{j})\text{ m/s}^2$. Its speed after $10\text{ seconds}$ is:
  \begin{multicols}{2}
  \begin{enumerate}[label=\textbf{\textcolor{primaryDark}{(\Alph*)}},itemsep=2pt]
    \item $7\sqrt{2}\text{ m/s}$
    \item $7\text{ m/s}$
    \item $8.5\text{ m/s}$
    \item $10\text{ m/s}$
  \end{enumerate}
  \end{multicols}

  % Q31
  \item The coordinates of a particle moving in the $xy$-plane at any instant of time $t$ are given by $x(t) = 4t^2$ and $y(t) = 3t^2$ (in meters). The speed of the particle at that instant is:
  \begin{multicols}{2}
  \begin{enumerate}[label=\textbf{\textcolor{primaryDark}{(\Alph*)}},itemsep=2pt]
    \item $10t\text{ m/s}$
    \item $5t\text{ m/s}$
    \item $7t\text{ m/s}$
    \item $25t\text{ m/s}$
  \end{enumerate}
  \end{multicols}

  % Q32
  \item A body is projected with velocity $v_1$ from point $A$ at an angle of $30^\circ$ with the horizontal. At the same instant, another body is projected vertically upwards from point $B$ (directly below the highest point of the first trajectory) with velocity $v_2$. For the two bodies to collide, the ratio $\dfrac{v_2}{v_1}$ must be:
  \begin{multicols}{2}
  \begin{enumerate}[label=\textbf{\textcolor{primaryDark}{(\Alph*)}},itemsep=2pt]
    \item $0.5$
    \item $1.0$
    \item $\dfrac{\sqrt{3}}{2}$
    \item $\dfrac{2}{\sqrt{3}}$
  \end{enumerate}
  \end{multicols}

  % Q33
  \item A rifle bullet is fired horizontally with a speed of $1500\text{ m/s}$ in order to hit a target $100\text{ m}$ away. Taking $g = 10\text{ m/s}^2$, the gun should be aimed:
  \begin{multicols}{2}
  \begin{enumerate}[label=\textbf{\textcolor{primaryDark}{(\Alph*)}},itemsep=2pt]
    \item $15\text{ cm}$ above the target
    \item $10\text{ cm}$ above the target
    \item $2.22\text{ cm}$ above the target
    \item Directly towards the target
  \end{enumerate}
  \end{multicols}

  % Q34
  \item Two projectiles $A$ and $B$ thrown with speeds in the ratio $1 : \sqrt{3}$ acquire the same maximum height. If $A$ is thrown at an angle of $30^\circ$ with the horizontal, the angle of projection of $B$ is:
  \begin{multicols}{2}
  \begin{enumerate}[label=\textbf{\textcolor{primaryDark}{(\Alph*)}},itemsep=2pt]
    \item $0^\circ$
    \item $\sin^{-1}(1/8)$
    \item $\sin^{-1}(1/\sqrt{12})$
    \item $\sin^{-1}(1/2)$
  \end{enumerate}
  \end{multicols}

  % Q35
  \item A ball is projected with a certain speed from the ground at a distance of $30\text{ m}$ from a vertical wall. If the angle of projection is $45^\circ$, the ball just clears the wall and strikes the ground $10\text{ m}$ past the wall. The height of the wall is:
  \begin{multicols}{2}
  \begin{enumerate}[label=\textbf{\textcolor{primaryDark}{(\Alph*)}},itemsep=2pt]
    \item $7.5\text{ m}$
    \item $10.0\text{ m}$
    \item $12.5\text{ m}$
    \item $15.0\text{ m}$
  \end{enumerate}
  \end{multicols}

  % Q36
  \item A ball rolls off the top of a staircase horizontally with velocity $u$. If the steps are $h$ meters high and $b$ meters wide, the ball will hit the edge of the $n^{\text{th}}$ step if:
  \begin{multicols}{2}
  \begin{enumerate}[label=\textbf{\textcolor{primaryDark}{(\Alph*)}},itemsep=2pt]
    \item $n = \dfrac{2hu^2}{gb}$
    \item $n = \dfrac{2hu}{gb^2}$
    \item $n = \dfrac{2hu^2}{gb^2}$
    \item $n = \dfrac{hu^2}{gb^2}$
  \end{enumerate}
  \end{multicols}

  % Q37
  \item A stone projected with velocity $u$ at angle $\theta$ reaches maximum height $H_1$. When projected with the same speed at angle $(90^\circ - \theta)$, it reaches height $H_2$. The relation between horizontal range $R$ and heights $H_1, H_2$ is:
  \begin{multicols}{2}
  \begin{enumerate}[label=\textbf{\textcolor{primaryDark}{(\Alph*)}},itemsep=2pt]
    \item $R = 4\sqrt{H_1 H_2}$
    \item $R = 4(H_1 - H_2)$
    \item $R = 4(H_1 + H_2)$
    \item $R = 2\sqrt{H_1 H_2}$
  \end{enumerate}
  \end{multicols}

  % Q38
  \item A body is projected with speed $u$ at an angle $\alpha$ with the horizontal. The speed of the body when its velocity vector makes an angle $\beta$ with the vertical is:
  \begin{multicols}{2}
  \begin{enumerate}[label=\textbf{\textcolor{primaryDark}{(\Alph*)}},itemsep=2pt]
    \item $\dfrac{u\cos\alpha}{\cos\beta}$
    \item $\dfrac{u\cos\alpha}{\sin\beta}$
    \item $\dfrac{u\sin\alpha}{\cos\beta}$
    \item $\dfrac{u\sin\alpha}{\sin\beta}$
  \end{enumerate}
  \end{multicols}

  % Q39
  \item An artillery piece with fixed muzzle speed has maximum range $R$. To hit a target at distance $R/2$ on the same horizontal level, the elevation angle of the gun should be:
  \begin{multicols}{2}
  \begin{enumerate}[label=\textbf{\textcolor{primaryDark}{(\Alph*)}},itemsep=2pt]
    \item $10^\circ$
    \item $18^\circ$
    \item $12^\circ$
    \item $15^\circ$
  \end{enumerate}
  \end{multicols}

  % Q40
  \item Building $P$ is $150\text{ m}$ tall and is at a horizontal distance of $100\text{ m}$ from building $Q$. A ball thrown horizontally from the roof of $P$ enters an open window in $Q$ located $25\text{ m}$ above the ground. The speed with which the ball was thrown is: (Take $g = 10\text{ m/s}^2$)
  \begin{multicols}{2}
  \begin{enumerate}[label=\textbf{\textcolor{primaryDark}{(\Alph*)}},itemsep=2pt]
    \item $20\text{ m/s}$
    \item $30\text{ m/s}$
    \item $40\text{ m/s}$
    \item $50\text{ m/s}$
  \end{enumerate}
  \end{multicols}

  % Q41
  \item Two bodies are projected horizontally in opposite directions with velocities $u_1$ and $u_2$ from the top of a building. If their velocity vectors become perpendicular to each other after falling through a vertical distance $h$, then $h$ is given by:
  \begin{multicols}{2}
  \begin{enumerate}[label=\textbf{\textcolor{primaryDark}{(\Alph*)}},itemsep=2pt]
    \item $h = \dfrac{u_1 u_2}{2g}$
    \item $h = \dfrac{2u_1 u_2}{g}$
    \item $h = \dfrac{u_1 u_2}{g}$
    \item $h = \dfrac{\sqrt{u_1 u_2}}{2g}$
  \end{enumerate}
  \end{multicols}

  % Q42
  \item A body is projected with velocity $u$ at an angle $\theta$ with the horizontal. The velocity of the body will become perpendicular to its initial velocity of projection after a time $t$ given by:
  \begin{multicols}{2}
  \begin{enumerate}[label=\textbf{\textcolor{primaryDark}{(\Alph*)}},itemsep=2pt]
    \item $\dfrac{2u\sin\theta}{g}$
    \item $\dfrac{u\sin\theta}{g}$
    \item $\dfrac{2u}{g\sin\theta}$
    \item $\dfrac{u}{g\sin\theta}$
  \end{enumerate}
  \end{multicols}

  % Q43
  \item A body is projected from the ground with speed $u$ at angle $\theta$ with the horizontal. At the highest point, its speed is $\sqrt{2/3}$ times its speed at half the maximum height. The angle of projection $\theta$ is:
  \begin{multicols}{2}
  \begin{enumerate}[label=\textbf{\textcolor{primaryDark}{(\Alph*)}},itemsep=2pt]
    \item $30^\circ$
    \item $45^\circ$
    \item $60^\circ$
    \item $\sin^{-1}(2/3)$
  \end{enumerate}
  \end{multicols}

  % Q44
  \item A body is projected from the ground with velocity $u$ at angle $\theta$. The average velocity of the body between the point of projection and the highest point of its trajectory is:
  \begin{multicols}{2}
  \begin{enumerate}[label=\textbf{\textcolor{primaryDark}{(\Alph*)}},itemsep=2pt]
    \item $\dfrac{u}{2}(1 + \cos\theta)$
    \item $\dfrac{u}{2}\sqrt{1 + \cos^2\theta}$
    \item $\dfrac{u}{2}\sqrt{1 + 2\cos^2\theta}$
    \item $\dfrac{u}{2}\sqrt{1 + 3\cos^2\theta}$
  \end{enumerate}
  \end{multicols}

  % Q45
  \item A string can withstand a maximum tension of $25\text{ N}$. What is the maximum speed with which a body of mass $1\text{ kg}$ can be whirled in a horizontal circle using a string of length $1\text{ m}$?
  \begin{multicols}{2}
  \begin{enumerate}[label=\textbf{\textcolor{primaryDark}{(\Alph*)}},itemsep=2pt]
    \item $2.5\text{ m/s}$
    \item $5.0\text{ m/s}$
    \item $7.5\text{ m/s}$
    \item $10.0\text{ m/s}$
  \end{enumerate}
  \end{multicols}

  % Q46
  \item A car travels at a speed of $36\text{ km/h}$ on a level road where the coefficient of static friction is $\mu_s = 0.8$. If $g = 10\text{ m/s}^2$, the car will skid while negotiating a curve if the radius of the turn is:
  \begin{multicols}{2}
  \begin{enumerate}[label=\textbf{\textcolor{primaryDark}{(\Alph*)}},itemsep=2pt]
    \item $10\text{ m}$
    \item $13\text{ m}$
    \item $14\text{ m}$
    \item $16\text{ m}$
  \end{enumerate}
  \end{multicols}

  % Q47
  \item A body of mass $0.5\text{ kg}$ is whirled in a vertical circle with an angular frequency of $10\text{ rad/s}$. If the radius of the circle is $0.5\text{ m}$, the tension in the string at the topmost point of the vertical circle is: (Take $g = 10\text{ m/s}^2$)
  \begin{multicols}{2}
  \begin{enumerate}[label=\textbf{\textcolor{primaryDark}{(\Alph*)}},itemsep=2pt]
    \item $10\text{ N}$
    \item $20\text{ N}$
    \item $30\text{ N}$
    \item $40\text{ N}$
  \end{enumerate}
  \end{multicols}

  % Q48
  \item The kinetic energy of a particle of mass $m$ moving in a circle of radius $r$ depends on distance covered $s$ as $K = as^2$, where $a$ is a positive constant. The total force acting on the particle as a function of $s$ is:
  \begin{multicols}{2}
  \begin{enumerate}[label=\textbf{\textcolor{primaryDark}{(\Alph*)}},itemsep=2pt]
    \item $2as\sqrt{1 + \left(\dfrac{s}{r}\right)^2}$
    \item $3as\sqrt{1 + \left(\dfrac{s}{r}\right)^2}$
    \item $as\sqrt{1 + \left(\dfrac{s}{r}\right)^2}$
    \item $4as\sqrt{1 + \left(\dfrac{s}{r}\right)^2}$
  \end{enumerate}
  \end{multicols}

  % Q49
  \item A coin of mass $10\text{ g}$ is placed at a distance of $30\text{ cm}$ from the center of a turn-table. The table rotates at $30\text{ rpm}$ about a vertical axis. The minimum coefficient of static friction to prevent the coin from sliding off is approximately: (Take $\pi^2 \approx 10, g = 10\text{ m/s}^2$)
  \begin{multicols}{2}
  \begin{enumerate}[label=\textbf{\textcolor{primaryDark}{(\Alph*)}},itemsep=2pt]
    \item $0.1$
    \item $0.2$
    \item $0.3$
    \item $0.4$
  \end{enumerate}
  \end{multicols}

  % Q50
  \item A conical pendulum has a string of length $L = 50\text{ cm}$ and a bob of mass $m = 200\text{ g}$. The bob revolves in a horizontal circle such that the string makes an angle $\theta = 60^\circ$ with the vertical. The tension in the string is: (Take $g = 10\text{ m/s}^2$)
  \begin{multicols}{2}
  \begin{enumerate}[label=\textbf{\textcolor{primaryDark}{(\Alph*)}},itemsep=2pt]
    \item $1\text{ N}$
    \item $2\text{ N}$
    \item $3\text{ N}$
    \item $4\text{ N}$
  \end{enumerate}
  \end{multicols}

\end{enumerate}

\newpage

% ==================== LEVEL 3 ====================
\begin{tcolorbox}[l3banner]
  {\Large \textbf{\textsf{LEVEL 3: JEE ADVANCED \& NUMERICAL CHALLENGE (25 Questions)}}}\\[2pt]
  {\footnotesize \textbf{Focus:} Incline Projections $\vert$ Turntable Kinematics $\vert$ Pursuit Paths $\vert$ Integer Numericals (Q70--Q75)}
\end{tcolorbox}
\vspace{6pt}

\begin{enumerate}[label=\textbf{\textcolor{l3Crimson}{Q\arabic*.}},leftmargin=24pt,itemsep=10pt,start=51]

  % Q51
  \item A particle is projected with a certain velocity at an angle $\alpha$ above the horizontal from the foot of an inclined plane of inclination $30^\circ$. If the particle strikes the plane normally, then $\alpha$ is equal to:
  \begin{multicols}{2}
  \begin{enumerate}[label=\textbf{\textcolor{primaryDark}{(\Alph*)}},itemsep=2pt]
    \item $30^\circ + \tan^{-1}\left(\dfrac{\sqrt{3}}{2}\right)$
    \item $45^\circ$
    \item $60^\circ$
    \item $30^\circ + \tan^{-1}(2\sqrt{3})$
  \end{enumerate}
  \end{multicols}

  % Q52
  \item A boat moves relative to river water with a velocity $v$ which is $n$ times less than the river flow velocity $u$ ($v = u/n$, $n > 1$). At what angle $\theta$ to the stream direction must the boat head to minimize drifting downstream?
  \begin{multicols}{2}
  \begin{enumerate}[label=\textbf{\textcolor{primaryDark}{(\Alph*)}},itemsep=2pt]
    \item $\dfrac{\pi}{2} + \sin^{-1}\left(\dfrac{1}{n}\right)$
    \item $\cos^{-1}\left(\dfrac{1}{n}\right)$
    \item $\tan^{-1}\left(\dfrac{1}{\sqrt{n^2 - 1}}\right)$
    \item $\dfrac{\pi}{2} - \sin^{-1}\left(\dfrac{1}{n}\right)$
  \end{enumerate}
  \end{multicols}

  % Q53
  \item Two cars $A$ and $B$ are moving west-to-east at $20\text{ m/s}$ and south-to-north at $15\text{ m/s}$ along perpendicular crossroads. Initially $A$ is $500\text{ m}$ west and $B$ is $400\text{ m}$ south of the intersection. The shortest distance between them during their subsequent motion is:
  \begin{multicols}{2}
  \begin{enumerate}[label=\textbf{\textcolor{primaryDark}{(\Alph*)}},itemsep=2pt]
    \item $10\text{ m}$
    \item $15\text{ m}$
    \item $20\text{ m}$
    \item $25\text{ m}$
  \end{enumerate}
  \end{multicols}

  % Q54
  \item An airplane is observed by two observers traveling at $60\text{ km/h}$ in vehicles moving in opposite directions along a straight road. To one observer, the airplane appears to cross the road track at right angles ($90^\circ$), while to the other it appears to cross at $45^\circ$. At what angle $\theta$ does the airplane actually cross the road track relative to the ground?
  \begin{multicols}{2}
  \begin{enumerate}[label=\textbf{\textcolor{primaryDark}{(\Alph*)}},itemsep=2pt]
    \item $\tan^{-1}(2)$
    \item $\tan^{-1}(1/2)$
    \item $\tan^{-1}(3/2)$
    \item $\tan^{-1}(2/3)$
  \end{enumerate}
  \end{multicols}

  % Q55
  \item A man wishes to cross a river of width $d = 500\text{ m}$. His rowing speed relative to water is $3\text{ km/h}$ and the river current is $2\text{ km/h}$. If the man's walking speed along the bank is $5\text{ km/h}$, in what direction $\theta$ (relative to the normal to the bank) should he row to reach the directly opposite point in the minimum possible time?
  \begin{multicols}{2}
  \begin{enumerate}[label=\textbf{\textcolor{primaryDark}{(\Alph*)}},itemsep=2pt]
    \item $\sin^{-1}(7/8)$
    \item $\sin^{-1}(8/7)$
    \item $\sin^{-1}(2/3)$
    \item $\sin^{-1}(3/7)$
  \end{enumerate}
  \end{multicols}

  % Q56
  \item A gun placed on a straight horizontal road is fired at an angle of $45^\circ$ to hit a car moving away from the gun with uniform speed $10\sqrt{2}\text{ m/s}$. The car is at a distance of $150\text{ m}$ from the gun at the instant of firing. The muzzle speed of the shell must be: (Take $g = 10\text{ m/s}^2$)
  \begin{multicols}{2}
  \begin{enumerate}[label=\textbf{\textcolor{primaryDark}{(\Alph*)}},itemsep=2pt]
    \item $20\text{ m/s}$
    \item $30\text{ m/s}$
    \item $40\text{ m/s}$
    \item $50\text{ m/s}$
  \end{enumerate}
  \end{multicols}

  % Q57
  \item From the top of a tower of height $40\text{ m}$, a ball is projected upwards with a speed of $20\text{ m/s}$ at an angle of $30^\circ$ to the horizontal. Taking $g = 10\text{ m/s}^2$, the total time after which the ball strikes the ground is:
  \begin{multicols}{2}
  \begin{enumerate}[label=\textbf{\textcolor{primaryDark}{(\Alph*)}},itemsep=2pt]
    \item $1\text{ s}$
    \item $2\text{ s}$
    \item $3\text{ s}$
    \item $4\text{ s}$
  \end{enumerate}
  \end{multicols}

  % Q58
  \item A uniform meter rod of mass $m = 1.5\text{ kg}$ and length $L = 1\text{ m}$ rests with its upper end against a smooth vertical wall and its lower end on a rough horizontal floor. The minimum coefficient of static friction $\mu_s$ between the floor and the rod so that it can lean at an angle of $30^\circ$ with the floor without slipping is:
  \begin{multicols}{2}
  \begin{enumerate}[label=\textbf{\textcolor{primaryDark}{(\Alph*)}},itemsep=2pt]
    \item $0.88$
    \item $0.98$
    \item $0.87$
    \item $0.67$
  \end{enumerate}
  \end{multicols}

  % Q59
  \item A man running along a straight horizontal road with uniform velocity $\vec{u} = u\hat{i}$ observes rain falling vertically down ($-\hat{j}$). If he doubles his running speed to $2u\hat{i}$, he observes the rain falling at an angle $\theta$ with the vertical. The true velocity of the rain with respect to the ground is:
  \begin{multicols}{2}
  \begin{enumerate}[label=\textbf{\textcolor{primaryDark}{(\Alph*)}},itemsep=2pt]
    \item $u\hat{i} - u\hat{j}$
    \item $u\hat{i} - \dfrac{u}{\tan\theta}\hat{j}$
    \item $2u\hat{i} + u\cot\theta\hat{j}$
    \item $u\hat{i} + u\sin\theta\hat{j}$
  \end{enumerate}
  \end{multicols}

  % Q60
  \item Three particles $A, B$, and $C$ are situated at the vertices of an equilateral triangle of side $l$ at $t = 0$. Each particle moves with constant speed $v$ such that $A$ always heads towards $B$, $B$ towards $C$, and $C$ towards $A$. The time when they meet at the centroid of the triangle is:
  \begin{multicols}{2}
  \begin{enumerate}[label=\textbf{\textcolor{primaryDark}{(\Alph*)}},itemsep=2pt]
    \item $\dfrac{2l}{3v}$
    \item $\dfrac{3l}{2v}$
    \item $\dfrac{l}{v}$
    \item $\dfrac{3l}{4v}$
  \end{enumerate}
  \end{multicols}

  % Q61
  \item A block of mass $m$ is placed on a rough horizontal turntable at distance $r$ from the axis. The coefficient of friction is $\mu$. If the turntable starts from rest with uniform angular acceleration $\alpha$, the time after which the block begins to slip is:
  \begin{multicols}{2}
  \begin{enumerate}[label=\textbf{\textcolor{primaryDark}{(\Alph*)}},itemsep=2pt]
    \item $\left[\left(\dfrac{\mu g}{\alpha^2 r}\right)^2 - \dfrac{1}{\alpha^2}\right]^{1/4}$
    \item $\left[\left(\dfrac{\mu g}{\alpha r}\right)^2 - \dfrac{1}{\alpha^2}\right]^{1/4}$
    \item $\left[\dfrac{\mu g}{\alpha^2 r} - \dfrac{1}{\alpha}\right]^{1/2}$
    \item $\left[\left(\dfrac{\mu g}{\alpha^2 r}\right)^2 + \dfrac{1}{\alpha^2}\right]^{1/4}$
  \end{enumerate}
  \end{multicols}

  % Q62
  \item A projectile is thrown upwards making an angle of $60^\circ$ with the horizontal with an initial speed of $150\text{ m/s}$. Taking $g = 10\text{ m/s}^2$, the time after which its inclination with the horizontal becomes $45^\circ$ is:
  \begin{multicols}{2}
  \begin{enumerate}[label=\textbf{\textcolor{primaryDark}{(\Alph*)}},itemsep=2pt]
    \item $15(\sqrt{3} - 1)\text{ s}$
    \item $15(\sqrt{3} + 1)\text{ s}$
    \item $7.5(\sqrt{3} - 1)\text{ s}$
    \item $7.5(\sqrt{3} + 1)\text{ s}$
  \end{enumerate}
  \end{multicols}

  % Q63
  \item A car travels east at $25\text{ km/h}$. To the driver of the car, a bus appears to move towards the north with a speed of $25\sqrt{3}\text{ km/h}$. The true speed and direction of motion of the bus are:
  \begin{multicols}{2}
  \begin{enumerate}[label=\textbf{\textcolor{primaryDark}{(\Alph*)}},itemsep=2pt]
    \item $50\text{ km/h}, 30^\circ\text{ East of North}$
    \item $50\text{ km/h}, 60^\circ\text{ North of East}$
    \item $25\text{ km/h}, 30^\circ\text{ East of North}$
    \item $25\text{ km/h}, 60^\circ\text{ North of East}$
  \end{enumerate}
  \end{multicols}

  % Q64
  \item A police van moving on a highway at $30\text{ km/h}$ fires a bullet at a suspect's car speeding away in the same direction at $192\text{ km/h}$. If the muzzle speed of the bullet is $150\text{ m/s}$, with what speed does the bullet strike the suspect's car?
  \begin{multicols}{2}
  \begin{enumerate}[label=\textbf{\textcolor{primaryDark}{(\Alph*)}},itemsep=2pt]
    \item $120\text{ m/s}$
    \item $90\text{ m/s}$
    \item $125\text{ m/s}$
    \item $105\text{ m/s}$
  \end{enumerate}
  \end{multicols}

  % Q65
  \item A body is projected vertically upwards with velocity $u$. After time $t$, another body is projected vertically upwards from the same point with velocity $v$ ($v < u$). For the two bodies to meet as soon as possible, the time interval $t$ must be:
  \begin{multicols}{2}
  \begin{enumerate}[label=\textbf{\textcolor{primaryDark}{(\Alph*)}},itemsep=2pt]
    \item $\dfrac{u - v + \sqrt{u^2 + v^2}}{g}$
    \item $\dfrac{u - v + \sqrt{u^2 - v^2}}{g}$
    \item $\dfrac{u + v + \sqrt{u^2 - v^2}}{g}$
    \item $\dfrac{u - v + \sqrt{u^2 - v^2}}{2g}$
  \end{enumerate}
  \end{multicols}

  % Q66
  \item A body is projected horizontally with speed $u$ from the top $A$ of an inclined plane of inclination $\theta = 30^\circ$ with the horizontal. It lands at point $P$ on the incline. The straight-line distance $AP$ along the incline is:
  \begin{multicols}{2}
  \begin{enumerate}[label=\textbf{\textcolor{primaryDark}{(\Alph*)}},itemsep=2pt]
    \item $\dfrac{u^2}{2g}$
    \item $\dfrac{\sqrt{3}u^2}{2g}$
    \item $\dfrac{2u^2}{3g}$
    \item $\dfrac{4u^2}{3g}$
  \end{enumerate}
  \end{multicols}

  % Q67
  \item A body is projected with speed $u$ at angle $\alpha$ with the horizontal. After what time $t$ will its instantaneous velocity vector become perpendicular to its initial velocity vector?
  \begin{multicols}{2}
  \begin{enumerate}[label=\textbf{\textcolor{primaryDark}{(\Alph*)}},itemsep=2pt]
    \item $\dfrac{2u\sin\alpha}{g}$
    \item $\dfrac{u}{g\sin\alpha}$
    \item $\dfrac{u\cos\alpha}{g}$
    \item $\dfrac{u}{g\cos\alpha}$
  \end{enumerate}
  \end{multicols}

  % Q68
  \item A light elastic spring of natural length $L = 50\text{ cm}$ and spring constant $k = 2\text{ N/m}$ has a small body of mass $m = 20\text{ g}$ attached to its end. The system revolves in a horizontal circle at $30\text{ rpm}$. Taking $\pi^2 \approx 10$, the elongation of the spring and the tension in it are:
  \begin{multicols}{2}
  \begin{enumerate}[label=\textbf{\textcolor{primaryDark}{(\Alph*)}},itemsep=2pt]
    \item $0.05\text{ m}, 0.1\text{ N}$
    \item $0.06\text{ m}, 0.2\text{ N}$
    \item $0.075\text{ m}, 0.2\text{ N}$
    \item $0.015\text{ m}, 0.3\text{ N}$
  \end{enumerate}
  \end{multicols}

  % Q69
  \item A ball is projected horizontally from the roof of a building. It strikes the ground after time $t$ with its velocity vector making an angle $\theta$ with the horizontal. The speed of horizontal projection was:
  \begin{multicols}{2}
  \begin{enumerate}[label=\textbf{\textcolor{primaryDark}{(\Alph*)}},itemsep=2pt]
    \item $gt\cot\theta$
    \item $gt\tan\theta$
    \item $gt\cos\theta$
    \item $gt\sin\theta$
  \end{enumerate}
  \end{multicols}

  % Q70
  \item A ball is thrown with a velocity whose horizontal component is $12\text{ m/s}$ from a point $15\text{ m}$ above the ground and $6\text{ m}$ away from a vertical wall $18.75\text{ m}$ high in such a way to just clear the wall. Find the time (in seconds) when the ball reaches the ground. (Take $g = 10\text{ m/s}^2$).\\
  \textit{[Numerical Value Question: Enter integer value]}

  % Q71
  \item A golfer hits a ball with velocity $u = 52\text{ m/s}$ at an angle $\alpha$ where $\tan\alpha = 5/12$. The time duration (in seconds) for which the ball remains at least $15\text{ m}$ above the ground is: (Take $g = 10\text{ m/s}^2$).\\
  \textit{[Numerical Value Question: Enter integer value]}

  % Q72
  \item Two cars $A$ and $B$ are moving west to east at $20\text{ m/s}$ and south to north at $15\text{ m/s}$ along crossroads. Initially $A$ is $500\text{ m}$ from the intersection and $B$ is $400\text{ m}$ from the intersection. If the shortest distance between them during motion is $d$, find the value of $d / 5$ (in meters).\\
  \textit{[Numerical Value Question: Enter integer value]}

  % Q73
  \item A boat is sent across (perpendicular) a river with a velocity of $8\text{ km/h}$. If the resultant velocity of the boat relative to the ground is $10\text{ km/h}$, find the velocity of the river flow (in km/h).\\
  \textit{[Numerical Value Question: Enter integer value]}

  % Q74
  \item A ball thrown by one player reaches another player in $2\text{ s}$. Find the maximum height (in meters) attained by the ball above the point of projection. (Take $g = 10\text{ m/s}^2$).\\
  \textit{[Numerical Value Question: Enter integer value]}

  % Q75
  \item Rain is falling vertically downwards with a speed of $4\text{ km/h}$. A girl walks on a straight horizontal road with a speed of $3\text{ km/h}$. The magnitude of the apparent velocity of rain with respect to the girl is (in km/h):\\
  \textit{[Numerical Value Question: Enter integer value]}

\end{enumerate}

\newpage

% ==================== MASTER ANSWER KEY ====================
\begin{tcolorbox}[sectionbanner]
  {\Large \textbf{\textsf{DAY 02: MASTER ANSWER KEY (Q1 -- Q75)}}}
\end{tcolorbox}
\vspace{6pt}

\begin{center}
\renewcommand{\arraystretch}{1.3}
\small
\begin{tabular}{|c|c||c|c||c|c||c|c||c|c|}
\hline
\rowcolor{tableDarkHeader}
\textcolor{white}{\textbf{Q.}} & \textcolor{white}{\textbf{Ans.}} & \textcolor{white}{\textbf{Q.}} & \textcolor{white}{\textbf{Ans.}} & \textcolor{white}{\textbf{Q.}} & \textcolor{white}{\textbf{Ans.}} & \textcolor{white}{\textbf{Q.}} & \textcolor{white}{\textbf{Ans.}} & \textcolor{white}{\textbf{Q.}} & \textcolor{white}{\textbf{Ans.}} \\
\hline\hline
\rowcolor{white}
1 & \textbf{\textcolor{l1Green}{B}} & 16 & \textbf{\textcolor{l1Green}{B}} & 31 & \textbf{\textcolor{l2Amber}{A}} & 46 & \textbf{\textcolor{l2Amber}{A}} & 61 & \textbf{\textcolor{l3Crimson}{A}} \\
\hline
\rowcolor{tableStripe}
2 & \textbf{\textcolor{l1Green}{C}} & 17 & \textbf{\textcolor{l1Green}{B}} & 32 & \textbf{\textcolor{l2Amber}{A}} & 47 & \textbf{\textcolor{l2Amber}{B}} & 62 & \textbf{\textcolor{l3Crimson}{C}} \\
\hline
\rowcolor{white}
3 & \textbf{\textcolor{l1Green}{C}} & 18 & \textbf{\textcolor{l1Green}{A}} & 33 & \textbf{\textcolor{l2Amber}{C}} & 48 & \textbf{\textcolor{l2Amber}{A}} & 63 & \textbf{\textcolor{l3Crimson}{B}} \\
\hline
\rowcolor{tableStripe}
4 & \textbf{\textcolor{l1Green}{A}} & 19 & \textbf{\textcolor{l1Green}{A}} & 34 & \textbf{\textcolor{l2Amber}{C}} & 49 & \textbf{\textcolor{l2Amber}{C}} & 64 & \textbf{\textcolor{l3Crimson}{D}} \\
\hline
\rowcolor{white}
5 & \textbf{\textcolor{l1Green}{D}} & 20 & \textbf{\textcolor{l1Green}{B}} & 35 & \textbf{\textcolor{l2Amber}{A}} & 50 & \textbf{\textcolor{l2Amber}{D}} & 65 & \textbf{\textcolor{l3Crimson}{B}} \\
\hline
\rowcolor{tableStripe}
6 & \textbf{\textcolor{l1Green}{B}} & 21 & \textbf{\textcolor{l1Green}{C}} & 36 & \textbf{\textcolor{l2Amber}{C}} & 51 & \textbf{\textcolor{l3Crimson}{A}} & 66 & \textbf{\textcolor{l3Crimson}{D}} \\
\hline
\rowcolor{white}
7 & \textbf{\textcolor{l1Green}{B}} & 22 & \textbf{\textcolor{l1Green}{A}} & 37 & \textbf{\textcolor{l2Amber}{A}} & 52 & \textbf{\textcolor{l3Crimson}{A}} & 67 & \textbf{\textcolor{l3Crimson}{B}} \\
\hline
\rowcolor{tableStripe}
8 & \textbf{\textcolor{l1Green}{B}} & 23 & \textbf{\textcolor{l1Green}{C}} & 38 & \textbf{\textcolor{l2Amber}{B}} & 53 & \textbf{\textcolor{l3Crimson}{C}} & 68 & \textbf{\textcolor{l3Crimson}{A}} \\
\hline
\rowcolor{white}
9 & \textbf{\textcolor{l1Green}{D}} & 24 & \textbf{\textcolor{l1Green}{A}} & 39 & \textbf{\textcolor{l2Amber}{D}} & 54 & \textbf{\textcolor{l3Crimson}{A}} & 69 & \textbf{\textcolor{l3Crimson}{A}} \\
\hline
\rowcolor{tableStripe}
10 & \textbf{\textcolor{l1Green}{B}} & 25 & \textbf{\textcolor{l1Green}{A}} & 40 & \textbf{\textcolor{l2Amber}{A}} & 55 & \textbf{\textcolor{l3Crimson}{D}} & 70 & \textbf{\textcolor{l3Crimson}{3}} \\
\hline
\rowcolor{white}
11 & \textbf{\textcolor{l1Green}{A}} & 26 & \textbf{\textcolor{l2Amber}{C}} & 41 & \textbf{\textcolor{l2Amber}{A}} & 56 & \textbf{\textcolor{l3Crimson}{D}} & 71 & \textbf{\textcolor{l3Crimson}{2}} \\
\hline
\rowcolor{tableStripe}
12 & \textbf{\textcolor{l1Green}{A}} & 27 & \textbf{\textcolor{l2Amber}{A}} & 42 & \textbf{\textcolor{l2Amber}{D}} & 57 & \textbf{\textcolor{l3Crimson}{D}} & 72 & \textbf{\textcolor{l3Crimson}{4}} \\
\hline
\rowcolor{white}
13 & \textbf{\textcolor{l1Green}{A}} & 28 & \textbf{\textcolor{l2Amber}{D}} & 43 & \textbf{\textcolor{l2Amber}{B}} & 58 & \textbf{\textcolor{l3Crimson}{C}} & 73 & \textbf{\textcolor{l3Crimson}{6}} \\
\hline
\rowcolor{tableStripe}
14 & \textbf{\textcolor{l1Green}{A}} & 29 & \textbf{\textcolor{l2Amber}{C}} & 44 & \textbf{\textcolor{l2Amber}{D}} & 59 & \textbf{\textcolor{l3Crimson}{B}} & 74 & \textbf{\textcolor{l3Crimson}{5}} \\
\hline
\rowcolor{white}
15 & \textbf{\textcolor{l1Green}{B}} & 30 & \textbf{\textcolor{l2Amber}{A}} & 45 & \textbf{\textcolor{l2Amber}{B}} & 60 & \textbf{\textcolor{l3Crimson}{A}} & 75 & \textbf{\textcolor{l3Crimson}{5}} \\
\hline
\end{tabular}
\end{center}

\vspace{10pt}
\begin{tcolorbox}[theorycard]
\footnotesize
\textbf{\textcolor{primaryDark}{Self-Evaluation \& JEE Benchmarking Guidelines:}}
\begin{itemize}[leftmargin=12pt,itemsep=2pt]
  \item \textbf{Marking Scheme:} $+4$ for correct, $-1$ for incorrect in MCQs, $0$ unattempted. $+4$ for correct, $0$ incorrect in Numericals. Max Marks: 300.
  \item \textbf{\textcolor{l1Green}{Target $> 240$ Marks:}} Exceptional 2D mastery \& quick spatial visualization (Top 500 AIR pace).
  \item \textbf{\textcolor{l2Amber}{Target $190 - 240$ Marks:}} Strong JEE Main performance (Top 2500 AIR pace). Review relative motion sign conventions.
  \item \textbf{\textcolor{l3Crimson}{Target $< 190$ Marks:}} Rework missed Level 2 and Level 3 problems in Block 5.
\end{itemize}
\end{tcolorbox}

\end{document}
